\documentclass[10pt,a4paper]{amsart}
\usepackage[margin=19mm]{geometry}
\usepackage{amsmath,amssymb,mathtools}
\usepackage[scr=rsfs,cal=euler]{mathalfa}
\usepackage{xfrac}
\usepackage[unicode,hidelinks,bookmarks=false]{hyperref}
\newcommand{\ie}{i.e.\ }
\newcommand{\cf}{cf.\ }
\newcommand{\resp}{respectively\ }
\newcommand{\Subsection}{\subsection}
\providecommand{\nobreakdash}{\nobreak\hbox{-}}
\setlength{\emergencystretch}{2em}
\multlinegap0pt
\usepackage{amssymb}
\usepackage{mathtools}
\newtheorem{prop}{Proposition}
\newtheorem{thm}{Theorem}
\newcommand{\N}{\mathbb N}
\newcommand{\R}{\mathbb R}
\title{Landau-level splitting by a circular $\delta$ interaction: boundary overlap and the hard-wall crossover}
\author{Masahiro Kaminaga}
\date{Source: arXiv:2509.08377v3 (CC BY 4.0)}
\begin{document}
\maketitle

\noindent\emph{Source and licence.} Landau-level splitting by a circular $\delta$ interaction: boundary overlap and the hard-wall crossover. Version arXiv:2509.08377v3 (CC BY 4.0), distributed by arXiv under the Creative Commons Attribution 4.0 International licence, http://creativecommons.org/licenses/by/4.0/, as recorded in the arXiv licence metadata for this exact version and shown on the abstract page https://arxiv.org/abs/2509.08377v3. Full licence notice: \texttt{licenses/arxiv-2509.08377-CC-BY-4.0.txt}.\par\smallskip

\noindent\textit{Editorial scope.} Counted unit: the article's thm thm:shift (source0.tex lines 860--883). The article's complete original proof of it is delivered with it (source0.tex lines 885--921, one \texttt{proof} environment of 37 lines). No mathematics is written, repaired or re-derived: the statement and the proof are the article's own lines, in the article's own notation, and no retained line is substituted or re-flowed.  Retained uncounted as the source-local context the proof needs: lines 529--533; lines 670--674; lines 736--756; lines 745--750; lines 788--791.  Nothing was omitted from any retained span, so the package carries no omission record.  No retained line contains a citation, so the wrapper carries no bibliography and no bibliography entry.  No retained line contains a figure environment, an image-inclusion command or drawing code, so no image is delivered and the package declares no asset dependency.  Editorial formatting only, in addition: an AMS \texttt{amsart} preamble that restates the article's own declarations and macros used by the retained lines, namely the environments \texttt{prop}, \texttt{thm} on separate counters, together with the standard packages those lines need.  The retained environments print in this document as Proposition 1, Theorem 1 in order of appearance, whereas the article prints the same items with the numbering of its own full document.  The complete native source of the article as distributed in the arXiv e-print for this exact version is bundled unchanged as \texttt{sources/184819/source0.tex}.
\begin{equation}\label{eq:detmatch}
\phi_m(a;E)\psi_m'(a;E)
-\phi_m'(a;E)\psi_m(a;E)
-\alpha\phi_m(a;E)\psi_m(a;E)=0.
\end{equation}

\begin{equation}\label{eq:c-exact}
c_{n,m}
=aB e^{-x}\frac{n!}{(n+m)!}x^m
\bigl[L_n^{(m)}(x)\bigr]^2.
\end{equation}

\begin{prop}\label{prop:c-asym}
For fixed $n$, $a$, and $B$, as $m\to\infty$,
\begin{equation}\label{eq:c-asym-1}
c_{n,m}
=aB e^{-x}\frac{n!}{(n+m)!}x^m
\binom{n+m}{n}^2
\left(1-\frac{2nx}{m}+O(m^{-2})\right).
\end{equation}
Equivalently,
\begin{equation}\label{eq:c-asym-2}
c_{n,m}
=\frac{aB e^{-x}}{n!}
\frac{x^m m^n}{m!}
\left[1+\frac{\kappa_n(x)}{m}+O(m^{-2})\right],
\end{equation}
where
\begin{equation}\label{eq:kappa}
\kappa_n(x)=\frac{n(n+1)}{2}-2nx.
\end{equation}
In particular, $c_{n,m}$ decreases faster than every exponential in $m$.
\end{prop}

\begin{equation}
c_{n,m}
=\frac{aB e^{-x}}{n!}
\frac{x^m m^n}{m!}
\left[1+\frac{\kappa_n(x)}{m}+O(m^{-2})\right],
\end{equation}

\begin{equation}\label{eq:regular-estimate}
r_{n,m}(E)=\frac{a}{2m}+\frac{ax}{2m^2}+O(m^{-3}),
\qquad x=\frac{Ba^2}{2},
\end{equation}

\begin{thm}\label{thm:shift}
Fix $n\in\N_0$, $B>0$, $a>0$, and $\alpha\in\R\setminus\{0\}$. For all
sufficiently large $m\ge0$, there is exactly one eigenvalue
$E_{n,m}^{(\alpha)}$ of the $m$-th radial channel in
$(\Lambda_n-\varepsilon,\Lambda_n+\varepsilon)$, where
$0<\varepsilon<B$ is fixed. It satisfies
\begin{equation}\label{eq:shift-c}
E_{n,m}^{(\alpha)}-\Lambda_n
=\alpha c_{n,m}
\left[1-\frac{\alpha a}{2m}+O(m^{-2})\right],
\qquad m\to\infty,
\end{equation}
where $c_{n,m}$ is the exact boundary weight in \eqref{eq:c-exact}. Consequently,
\begin{eqnarray}
E_{n,m}^{(\alpha)}-\Lambda_n
&=&\frac{\alpha aB e^{-x}}{n!}
\frac{x^m m^n}{m!}
\left[1+\frac{\kappa_n(x)-\alpha a/2}{m}
+O(m^{-2})\right],\label{eq:shift-explicit}
\\[-2mm]
&&\hspace{20mm} x=\frac{Ba^2}{2},\qquad
\kappa_n(x)=\frac{n(n+1)}2-2nx.\nonumber
\end{eqnarray}
\end{thm}

\begin{proof}
Put $c=c_{n,m}>0$ for large $m$ and write $E=\Lambda_n+\delta$.
Multiplication of the scalar equation by $\delta$ gives
\begin{equation}\label{eq:Fdelta}
F_m(\delta):=\delta[1+\alpha r_{n,m}(\Lambda_n+\delta)]-\alpha c=0.
\end{equation}
The regular-part lemma gives $|\alpha r_{n,m}|<1/4$ throughout the
fixed neighborhood once $m$ is large. The values of $F_m$ at $\delta=\alpha c/2$ and
$\delta=2\alpha c$ have opposite signs, irrespective of the sign of
$\alpha$.
Both points belong to that neighborhood by Proposition~\ref{prop:c-asym}.
Thus there is a real root with the sign of $\alpha$ and
$|\delta|\le2|\alpha|c$. Moreover,
$$
F_m'(\delta)=1+\alpha r_{n,m}(\Lambda_n+\delta)
+\alpha\delta r_{n,m}'(\Lambda_n+\delta)>0
$$
throughout the neighborhood for large $m$. Hence this root is unique.
The center $\delta=0$ is not an eigenvalue: its free eigenfunction has
nonzero trace, and \eqref{eq:detmatch} at that energy is nonzero.
Any other eigenvalue in the neighborhood would satisfy \eqref{eq:Fdelta}.
The bound on $|\delta|$ also proves convergence to $\Lambda_n$.

At the root we have
\begin{equation}\label{eq:delta-solve}
\delta=\frac{\alpha c_{n,m}}
{1+\alpha r_{n,m}(\Lambda_n+\delta)}.
\end{equation}
Inserting \eqref{eq:regular-estimate} and expanding the denominator proves
\eqref{eq:shift-c}. Together with \eqref{eq:c-asym-2}, this gives
\eqref{eq:shift-explicit}. In fact the same proof gives the stronger form
\begin{equation}\label{eq:shift-second}
\frac{E_{n,m}^{(\alpha)}-\Lambda_n}{\alpha c_{n,m}}
=1-\frac{\alpha a}{2m}
+\frac{(\alpha a/2)^2-\alpha ax/2}{m^2}+O(m^{-3}).
\end{equation}
\end{proof}

\end{document}
